% LaTeX template for Finance summaries
% Stu Alden
%----------------------------------------------------------------------------------------
% Bonds and Stocks
\documentclass[14pt]{extarticle}

\usepackage{conceptsheet}

%---------------------------------------
\begin{document}

\finmathtitle{6. Bonds \& Stocks}


We can analyze {stocks} and {bonds} using all of the compound interest techniques we've
learned so far, by examing the expected cashflows.

\begin{itemize}
  \item A \textbf{bond} (purchased at issue, or later on the seconday market) can be viewed as a combination of
    \begin{itemize}
      \item An \underline{annuity} of coupon payments, plus
      \item A \underline{single} payment of the face amount at maturity
    \end{itemize}
  \item A \textbf{stock} can be viewed as
    \begin{itemize}
      \item An annuity of dividends (if applicable), perhaps varying over time, plus
      \item The proceeds received from the sale of the stock at some future date.
    \end{itemize}
\end{itemize}
\vspace{.1in}

Bond cashflows are generally far more predictable than stock cashflows, but the basic principles
for analysis are the same.

In most cases we either know a desired a yield rate (i) and we are trying to solve for A (the
price, namely the present value of the cashflows at the time of purchase) or we know the
purchase price and we want to estimate the yield.  Either way, we can use compound interest and
the annuity formulas.  So, if

\begin{description}
  \item \textbf{R} - Amount of each payment (coupon or dividend)
  \item\textbf{n} - Number of periods until bond maturity or sale of stock
  \item\textbf{M} - Maturity value of bond / sale price of stock
  \item\textbf{i} - Compound interest rate
  \item\textbf{A} - Present Value (purchase price) of the cashflows
\end{description}
\vspace{.1in}

Then


\begin{align*}
  A & = R \cdot \frac{1 - v^n}{i} + M \cdot v^n
\end{align*}

%    & = R \cdot \frac{1 - v^n}{i}

% \begin{align*}
%     S & = R \cdot \frac{(1 + i)^n-1}{i} \\
% \end{align*}
% \vspace{-.2in}
\begin{flushleft}
  where, as always, {\large $ v = \frac{1}{(1+i)} = (1+i)^{-1} $ }.
\end{flushleft}
% \vspace{-.1in}
\vspace{.2in}

% \begin{center}
%     \begin{tikzpicture}
%         %draw horizontal line with squiggle in middle
%         \draw (0,0) -- (5,0);
%         \draw[decorate,decoration={snake,pre length=5mm, post length=5mm}] (5,0) -- (7,0);
%         \draw (7,0) -- (10,0);

%         %draw vertical tick-marks
%         \foreach \x in {0,2,4,8,10}
%         \draw (\x cm,3pt) -- (\x cm,-3pt);

%         %Label "nodes"
%         \draw (0,0) node[below=3pt] {$ 0 $} node[above=3pt] {$  $};
%         \draw (0,0) node[below=18pt] {$ A $} node[above=3pt] {$  $};

%         \draw (2,0) node[below=3pt] {$ 1 $} node[above=3pt] {$ R $};
%         \draw (4,0) node[below=3pt] {$ 2 $} node[above=3pt] {$ R $};
%         \draw (6,0) node[below=3pt] {$  $} node[above=3pt] {$  $};
%         \draw (8,0) node[below=3pt] {$ n - 1 $} node[above=3pt] {$ R $};
%         \draw (10,0) node[below=3pt] {$ n $} node[above=3pt] {$ R $};
%         \draw (10,0) node[below=18pt] {$ S $} node[above=3pt] {$ R $};
%     \end{tikzpicture}
% \end{center}
% %\vspace{.1in}


\textbf{Notes:}

\begin{itemize}
  \item The price calculation (A) is generally straightforward if the yield (i) is known.
  \item The yield calculation, on the other hand, is generally not feasible with "analytic" (direct)
    methods, but is easy to estimate to any required precision with the aid of a computer.
\end{itemize}

\end{document}
